The lower the cutoff frequency, the more low-frequency information is preserved in the hybrid image. \\
This means that the resulting hybrid image will have more of the original image's overall structure and less of the high-frequency details. \\
Conversely, a higher cutoff frequency will allow more high-frequency information to pass through, resulting in a hybrid image that retains more of the fine details from the original images.\\
For example, if the cutoff frequency is set to a low value, the hybrid image will appear more blurred and less detailed, while a higher cutoff frequency will produce a sharper and more detailed hybrid image.\\
In this case, if the cutoff frequency is set to 2 (low value), the image will look more like the dog. \\
The image will have less sharper details but will have more of the structure of the orignal dog image. \\
If the cutoff frequency is set to a higher value of 20, the resulting hybrid image will have more of the finer sharp details. \\
The resulting image will look more like the cat image with a lot of the sharp details from the original cat image. \\